%! Solving Quadratics by Factoring — Unit 2, Lesson 2.2 — Student Packet Cover Sheet
% SLIDES-FIRST shape (2026-09-25; 2.2 is the pilot). 12pt; \coverbanner measures the title block.
% TWO scored rows: the slide handout (the student's notes, printed separately as
% lesson02_slides.pdf) and the homework — a printed DeltaMath set for this lesson (user choice).
% Homework is due the first class after two study halls — never "next class".
\documentclass[12pt]{article}
\usepackage{algebra2-article}
\usepackage{algebra2-boxes}

\begin{document}

\coverbanner{Unit 2 \quad Quadratic Functions}{Lesson 2.2 \quad Solving Quadratics by Factoring}

\namedateperiod
\vspace{0.18in}

\boxguard[14]
\begin{learningtargetbox}
\small
\begin{enumerate}[leftmargin=1.5em, itemsep=5pt]
  \item \ldots{} \textbf{factor} $x^2+bx+c$ by finding two integers with the right
        \textbf{product} and \textbf{sum}, taking out the \textbf{greatest common factor} first.
  \item \ldots{} factor $ax^2+bx+c$ by \textbf{grouping}.
  \item \ldots{} recognize and factor a \textbf{difference of squares} and a
        \textbf{perfect-square trinomial}, and know that a sum of squares does not factor.
  \item \ldots{} solve a quadratic equation with the \textbf{Zero Product Property}, and read its
        \textbf{roots} as the \textbf{$x$-intercepts} of its graph.
\end{enumerate}
\end{learningtargetbox}

\vspace{0.14in}

\boxguard[14]
\begin{tocbox}
\small
\renewcommand{\arraystretch}{1.5}
\begin{tabularx}{\linewidth}{@{} c l X r @{}}
\rowcolor{forestbg}
\textbf{\#} & \textbf{Component} & \textbf{Description} & \textbf{Score} \\
\midrule
1 & Slide Notes & The printed slides: the warm-up, four definitions with worked examples, and
                  every \emph{Now you try} worked in the notes column            & \blank{1.2cm} \\
2 & Homework (DeltaMath) & Factoring and solving by factoring ---
                  \textbf{scored; due the first class after two study halls}      & \blank{1.2cm} \\
\midrule
  & \multicolumn{2}{r}{\textbf{Total}} & \blank{1.2cm} \\
\end{tabularx}
\end{tocbox}

\vspace{0.14in}

\boxguard[8]
\begin{remindbox}
\small
To \textbf{factor} is to write a polynomial as a product --- take out the \textbf{GCF} first. For
$x^2+bx+c$, find two integers with product $c$ and sum $b$; for $ax^2+bx+c$, find two with
product $a\cdot c$ and sum $b$, split the middle term, and \textbf{group}. On sight:
$a^2-b^2=(a+b)(a-b)$ and $a^2\pm2ab+b^2=(a\pm b)^2$; a sum of squares does not factor. The
\textbf{Zero Product Property} says that if $AB=0$ then $A=0$ or $B=0$ --- so get everything on
one side \emph{equal to zero} before you factor, and never divide both sides by $x$ (that loses
the root $x=0$). The \textbf{roots} are the graph's $x$-intercepts.
\end{remindbox}

\end{document}
